% =====================================================================
%  related_work.tex  --  PHASE I, Task 1
%
%  Structured four-theme literature review + cross-comparison table
%  + explicit research-gap statement.
%
%  Drop-in usage (IEEEtran or elsarticle):
%      % =====================================================================
%  related_work.tex  --  PHASE I, Task 1
%
%  Structured four-theme literature review + cross-comparison table
%  + explicit research-gap statement.
%
%  Drop-in usage (IEEEtran or elsarticle):
%      % =====================================================================
%  related_work.tex  --  PHASE I, Task 1
%
%  Structured four-theme literature review + cross-comparison table
%  + explicit research-gap statement.
%
%  Drop-in usage (IEEEtran or elsarticle):
%      % =====================================================================
%  related_work.tex  --  PHASE I, Task 1
%
%  Structured four-theme literature review + cross-comparison table
%  + explicit research-gap statement.
%
%  Drop-in usage (IEEEtran or elsarticle):
%      \input{related_work}
%  Requires in the preamble:
%      \usepackage{booktabs}
%      \usepackage{threeparttable}   % for the table notes
%      \usepackage{multirow}
%      \usepackage{wasysym}          % \CIRCLE \LEFTcircle \Circle
%  Note: wasysym redefines \Box/\diamond; if that clashes with your
%  template, load it as \usepackage[nointegrals]{wasysym} or swap the
%  harvey-ball glyphs for {\ding{108}}/{\ding{109}} from pifont.
%  Bibliography: paper/references.bib
% =====================================================================

\section{Related Work}
\label{sec:related}

The problem addressed in this paper sits at the intersection of four
research communities that have so far developed largely in isolation:
(i)~unsupervised identification of machine operating states from
electrical measurements, (ii)~non-intrusive load monitoring (NILM),
(iii)~industrial energy optimisation at the process and system level,
and (iv)~predictive shutdown and decision support for manufacturing
equipment. We review each in turn, then summarise the resulting gap.


% ---------------------------------------------------------------------
\subsection{Machine State Identification}
\label{sec:rw:states}

The recognition that an electrically driven machine occupies a small
number of discrete operating regimes---typically \textsc{off},
\textsc{standby} (energised but not producing), and \textsc{working}---is
the foundation of essentially all energy-state analysis. ISO~14955-1
\cite{iso14955} formalises this decomposition as a design-time
requirement for energy-efficient machine tools, defining standby as the
operating mode in which a machine is ready for production but performs
no value-adding function. This definition is the one adopted throughout
the present work.

\subsubsection{Supervised approaches}
Where per-timestamp annotations are available, state identification
reduces to a standard classification problem, and support-vector
machines, random forests, and convolutional classifiers all perform
well. The practical obstacle in industrial deployments is that such
annotations essentially never exist: obtaining them requires either
instrumenting the machine's contactor and PLC directly or having an
operator observe each machine over an extended period. Both are
prohibitively expensive at facility scale, and neither transfers to a
new machine without repeating the effort. Supervised methods are
therefore poorly matched to the retrofit scenario that motivates this
paper.

\subsubsection{Unsupervised approaches}
Unsupervised state inference avoids the annotation cost entirely.
$k$-means \cite{lloyd1982least} is the simplest option but imposes
isotropic, equal-variance clusters, which is a poor model for
electrical states whose dispersion grows with load. Gaussian mixture
models (GMM) \cite{mclachlan2000finite}, fitted by
expectation--maximisation \cite{dempster1977em}, relax both assumptions:
each state receives its own mean and full covariance, and assignments
are soft, which matters precisely at the \textsc{standby}/\textsc{working}
boundary where the two regimes overlap. Density-based methods such as
DBSCAN \cite{ester1996dbscan} and HDBSCAN \cite{campello2015hdbscan}
make no parametric shape assumption and identify noise points
explicitly, while spectral clustering \cite{ng2001spectral} operates on
the affinity graph and can recover non-convex structure. Model order is
conventionally selected by the Bayesian information criterion
\cite{schwarz1978bic} or by silhouette analysis
\cite{rousseeuw1987silhouettes}.

Neither criterion selects a state count on records of the size
considered here, and this is worth stating because it determines the
design of the method proposed below. The BIC penalty grows as
$\tfrac{1}{2}p\log n$ while the likelihood gained by splitting a
non-Gaussian component grows linearly in $n$, so at $n>10^{6}$ samples
per machine the penalty is negligible and a mixture model keeps buying
components in order to approximate a distribution that is not a mixture
of $k$ Gaussians for any $k$. Scanning $k\in[2,8]$ over the eight
machines of Section~\ref{sec:multimachine}, BIC is still falling at the
boundary on seven of them, and the improvement from $k=4$ to $k=8$ is
$3.1$--$33.1\%$ of the criterion's own value. Silhouette behaves
similarly. Statistical order selection must therefore be supplemented
by an external, physical admissibility criterion rather than merely
confirmed by one --- which is the role of the validation battery of
Section~\ref{sec:method:accept_reject}, and the reason it is part of
the method here rather than an appendix to it.

\subsubsection{Temporal structure}
A limitation common to all of the above is that they treat samples as
independent and identically distributed. Applied to a 1\,Hz electrical
stream, an i.i.d.\ clustering assigns each second in isolation and
consequently produces rapid, physically impossible state alternation
around cluster boundaries---a failure mode we quantify as the
\emph{flicker rate} in Section~\ref{sec:results}. Hidden Markov models
\cite{rabiner1989tutorial} address this by making the state sequence
itself the object of inference: the transition matrix encodes state
persistence, and Viterbi decoding returns the globally most probable
sequence rather than a concatenation of independent decisions. Fox
\textit{et al.} \cite{fox2011sticky} show that an explicit bias toward
self-transition (a ``sticky'' prior) reduces over-segmentation when the
emission distributions overlap, which is the regime encountered here.

It is worth stating precisely what this machinery can and cannot do,
because the point is routinely overstated. In a Viterbi decode the cost
of changing state is $\log(p_{\text{self}}/p_{\text{switch}})$, which is
a few nats even for extremely persistent transition matrices and is
bounded above however strong the prior is made. The emission term
carries no such bound: full-covariance Gaussians fitted to low-noise,
multi-channel electrical measurements can differ by hundreds of nats
between adjacent states. Where that holds, the transition model is
arithmetically unable to influence the decoded path and the HMM
degenerates toward per-sample assignment. Guaranteeing a minimum state
duration therefore requires modelling duration explicitly---a hidden
semi-Markov formulation \cite{yu2010hsmm}, or an equivalent
duration constraint imposed at decode time---rather than a heavier
transition prior. Section~\ref{sec:results} quantifies this for the
present data.

\subsubsection{Threshold-free decision rules}
Where a scalar decision boundary is unavoidable---for instance when
converting a continuous power-factor distribution into a binary
load/no-load verdict---Otsu's criterion \cite{otsu1979threshold}
selects the threshold that maximises between-class variance, removing
the arbitrariness of a hand-tuned cut-off. Agreement among several
independently derived partitions can then be aggregated using cluster
ensemble methods \cite{strehl2002cluster}.


% ---------------------------------------------------------------------
\subsection{Non-Intrusive Load Monitoring}
\label{sec:rw:nilm}

NILM originates with Hart \cite{hart1992nonintrusive}, who showed that
individual appliance activity can be recovered from an aggregate meter
by detecting steady-state step changes in real and reactive power.
Subsequent surveys \cite{zeifman2011nilm,schirmer2023nilmreview} trace
the field's progression from event-based signature matching to
probabilistic and, more recently, deep-learning formulations. Kolter
and Jaakkola \cite{kolter2012afamap} cast disaggregation as inference in
an additive factorial HMM, making explicit the connection between NILM
and the temporal state models of Section~\ref{sec:rw:states}. The deep
learning era is represented by Neural NILM \cite{kelly2015neuralnilm},
sequence-to-point regression \cite{zhang2018seq2point}, and denoising
autoencoders \cite{bonfigli2018denoising}. Standardised toolkits
\cite{batra2014nilmtk} and evaluation protocols
\cite{makonin2015performance} have made results across this line of
work broadly comparable.

Two properties of this literature are decisive for the present paper.
First, it is overwhelmingly \emph{residential}. The canonical benchmarks
---REDD \cite{kolter2011redd} and UK-DALE---consist of household
appliances whose signatures are sparse, short, and well separated.
Industrial machinery behaves very differently: loads are large,
continuously modulated by process throughput, dominated by induction
motors whose reactive power draw persists at zero mechanical output, and
governed by shift schedules rather than occupant behaviour. Holmegaard
and Kj\ae{}rgaard \cite{holmegaard2016industrial} evaluate established
NILM algorithms on industrial loads and report substantially degraded
performance, attributing it to precisely these differences. Martins
\textit{et al.} \cite{martins2018imdeld} apply a deep generative model
to industrial disaggregation and, in doing so, release the IMDELD
dataset \cite{imdeld2018dataset}---eleven meters recording eight heavy
machines in a Brazilian poultry-feed pellet plant at 1\,Hz across a
five-month window, of which the covered fraction differs by machine
(Section~\ref{sec:method:data}). IMDELD is the benchmark used throughout this
paper; its authors characterise the machines as three-state systems but
supply no per-timestamp state annotation, which is the reason the
present work must proceed without labels and validate its partition
against physics instead.

Second, and more fundamentally, the objective of NILM is
\emph{attribution}: given an aggregate signal, determine which appliance
consumed what. It is a retrospective accounting exercise. NILM does not
forecast, and it does not act. No NILM system, industrial or
residential, closes the loop from an inferred state to an actuation
decision. That step belongs to the two literatures reviewed next.


% ---------------------------------------------------------------------
\subsection{Industrial Energy Optimisation}
\label{sec:rw:industrial}

The manufacturing engineering community has independently established
that non-productive energy dominates the energy budget of discrete-part
production. Gutowski \textit{et al.} \cite{gutowski2006electrical} and
Dahmus and Gutowski \cite{dahmus2004environmental} show that the
load-independent baseline of a machine tool---spindle idling, hydraulics,
coolant pumps, control electronics---frequently exceeds the energy
actually delivered to the workpiece, and that this fraction rises as
automation increases. Kara and Li \cite{kara2011unit} formalise the
relationship as a unit-process model in which specific energy
consumption is inversely related to material removal rate, so that a
lightly loaded or idle machine is intrinsically inefficient. Devoldere
\textit{et al.} \cite{devoldere2007improvement} quantify the resulting
improvement potential and identify idle-mode elimination as the single
largest lever available without redesigning the machine. Duflou
\textit{et al.} \cite{duflou2012towards} synthesise this body of work
into a systems-level roadmap, placing operational measures---switching
equipment off when it is not needed---at the machine and cell levels of
their hierarchy.

On the management side, ISO~50001 \cite{iso50001} provides the
plan--do--check--act framework within which energy performance
indicators are defined and tracked, and ISO~14955-1 \cite{iso14955}
supplies the operating-mode taxonomy that makes standby energy a
measurable quantity in the first place.

The limitation of this literature is methodological rather than
conceptual. Its energy models are typically built from controlled
laboratory measurements of a single machine under prescribed process
parameters, and its recommendations are design guidelines and audit
procedures. It does not address how a facility, given only the
electrical measurements it already collects and no process-parameter
instrumentation, should determine \emph{which} of its machines are idle,
\emph{when}, and for \emph{how long}. That inference problem is left
open.


% ---------------------------------------------------------------------
\subsection{Predictive Shutdown and Decision Support}
\label{sec:rw:decision}

The operations research community has studied the shutdown decision
itself in considerable depth. Mouzon \textit{et al.}
\cite{mouzon2007operational} establish the central trade-off: turning a
machine off saves its idle power but incurs a restart energy penalty and
a restart delay, so shutdown is economical only when the idle period
exceeds a break-even duration determined by the ratio of restart energy
to idle power. They propose dispatching rules and a mathematical
programming formulation that exploit this, and extend the treatment to
joint energy--tardiness objectives \cite{mouzon2008framework}. Chen
\textit{et al.} \cite{chen2013schedule} develop schedule-based operation
policies for production lines, Sun and Li \cite{sun2013opportunity}
estimate real-time control opportunities from buffer state, Li and Sun
\cite{li2016energy} formulate dynamic energy control as a Markov
decision process, Frigerio and Matta \cite{frigerio2015energy} derive
control strategies for machine tools under stochastic job arrivals, and
He \textit{et al.} \cite{he2015energy} embed energy responsiveness in
job-shop machine selection and sequencing.

This literature is mathematically mature, and the break-even
formulation of Mouzon \textit{et al.} is the natural objective for the
decision layer of any standby-elimination system. Its universal
assumption, however, is that the idle interval is \emph{known}---either
because it is dictated by a deterministic production schedule, or
because it is drawn from an analytically tractable arrival process
calibrated by the modeller. The policies are evaluated in simulation or
on analytical queueing models. What is absent is any mechanism for
obtaining the idle-duration input from a machine's own measured
electrical behaviour, which is the only signal available in a retrofit
installation. In other words: this community can decide optimally given
a duration, and the NILM community can infer a state from
measurements, but no work connects the two.


% ---------------------------------------------------------------------
\subsection{Synthesis and Research Gap}
\label{sec:rw:gap}

Table~\ref{tab:related} compares representative works from the four
themes against the capabilities required for an autonomous
standby-elimination system: whether the method infers machine state
without labels, whether it models temporal structure explicitly, whether
it forecasts the \emph{duration} of a future non-productive interval,
whether it converts that forecast into an economically justified control
decision, and whether it is validated on real industrial measurements.

\begin{table*}[!t]
\centering
\caption{Comparison of representative prior work against the
capabilities required for autonomous standby elimination.
\CIRCLE~= provided, \LEFTcircle~= partially provided, \Circle~= not
provided, --~= not applicable.}
\label{tab:related}
\begin{threeparttable}
\small
\setlength{\tabcolsep}{4pt}
\begin{tabular}{@{}lllccccc@{}}
\toprule
\multirow{2}{*}{\textbf{Work}} &
\multirow{2}{*}{\textbf{Method family}} &
\multirow{2}{*}{\textbf{Data domain}} &
\textbf{Labels} &
\textbf{Temporal} &
\textbf{Duration} &
\textbf{Economic} &
\textbf{Real ind.} \\
 & & & \textbf{required} & \textbf{model} & \textbf{forecast} &
 \textbf{decision} & \textbf{data} \\
\midrule
\multicolumn{8}{@{}l}{\textit{Theme A --- machine state identification}}\\
Lloyd \cite{lloyd1982least}          & $k$-means                    & generic                & no  & \Circle    & \Circle & \Circle & \Circle \\
McLachlan \& Peel \cite{mclachlan2000finite} & Gaussian mixture     & generic                & no  & \Circle    & \Circle & \Circle & \Circle \\
Ester \textit{et al.} \cite{ester1996dbscan} & DBSCAN               & generic                & no  & \Circle    & \Circle & \Circle & \Circle \\
Fox \textit{et al.} \cite{fox2011sticky}     & sticky HDP-HMM       & audio                  & no  & \CIRCLE    & \Circle & \Circle & \Circle \\
\addlinespace
\multicolumn{8}{@{}l}{\textit{Theme B --- non-intrusive load monitoring}}\\
Hart \cite{hart1992nonintrusive}     & event/steady-state signature & residential            & yes\tnote{a} & \LEFTcircle & \Circle & \Circle & \Circle \\
Kolter \& Jaakkola \cite{kolter2012afamap} & additive factorial HMM & REDD (residential)     & no  & \CIRCLE    & \Circle & \Circle & \Circle \\
Kelly \& Knottenbelt \cite{kelly2015neuralnilm} & deep NN (dAE, LSTM) & UK-DALE (residential) & yes & \LEFTcircle & \Circle & \Circle & \Circle \\
Zhang \textit{et al.} \cite{zhang2018seq2point} & seq2point CNN     & UK-DALE, REDD          & yes & \LEFTcircle & \Circle & \Circle & \Circle \\
Holmegaard \& Kj\ae{}rgaard \cite{holmegaard2016industrial} & NILM algorithm evaluation & industrial & yes & \LEFTcircle & \Circle & \Circle & \CIRCLE \\
Martins \textit{et al.} \cite{martins2018imdeld} & deep generative disaggregation & IMDELD (industrial) & yes & \LEFTcircle & \Circle & \Circle & \CIRCLE \\
\addlinespace
\multicolumn{8}{@{}l}{\textit{Theme C --- industrial energy optimisation}}\\
Gutowski \textit{et al.} \cite{gutowski2006electrical} & thermodynamic/empirical process model & lab machine tools & -- & \Circle & \Circle & \Circle & \LEFTcircle \\
Kara \& Li \cite{kara2011unit}       & unit-process energy model    & lab machine tools      & --  & \Circle    & \Circle & \Circle & \LEFTcircle \\
Duflou \textit{et al.} \cite{duflou2012towards} & systems-level roadmap & survey             & --  & \Circle    & \Circle & \LEFTcircle & \LEFTcircle \\
\addlinespace
\multicolumn{8}{@{}l}{\textit{Theme D --- predictive shutdown and decision support}}\\
Mouzon \textit{et al.} \cite{mouzon2007operational} & dispatching rules, MILP & simulated single machine & -- & \Circle & \Circle & \CIRCLE & \Circle \\
Chen \textit{et al.} \cite{chen2013schedule} & schedule-based control  & simulated line     & --  & \LEFTcircle & \Circle & \CIRCLE & \Circle \\
Frigerio \& Matta \cite{frigerio2015energy} & stochastic control policy & analytical model  & --  & \CIRCLE    & \LEFTcircle\tnote{b} & \CIRCLE & \Circle \\
He \textit{et al.} \cite{he2015energy} & energy-responsive scheduling & job-shop model      & --  & \Circle    & \Circle & \CIRCLE & \Circle \\
\midrule
\textbf{This work}                   & \textbf{physics-gated GMM--HMM $\rightarrow$ duration forecast $\rightarrow$ optimisation} & \textbf{IMDELD (industrial)} & \textbf{no} & \CIRCLE & \CIRCLE & \CIRCLE & \CIRCLE \\
\bottomrule
\end{tabular}
\begin{tablenotes}[flushleft]\footnotesize
\item[a] Hart's formulation is unsupervised in its clustering step but
requires a manually curated appliance signature library to name the
recovered clusters.
\item[b] Idle duration is modelled as a known stochastic arrival
process calibrated by the analyst, not forecast from measured machine
behaviour.
\end{tablenotes}
\end{threeparttable}
\end{table*}

Reading Table~\ref{tab:related} column-wise makes the gap explicit.
Every capability required for autonomous standby elimination has been
solved somewhere, but the last column of any single row is always
empty:

\begin{itemize}
  \item \textbf{Label-free state inference with temporal structure}
  exists (Theme~A, and \cite{kolter2012afamap} in Theme~B), but is
  demonstrated on generic or residential data and terminates at the
  labelled sequence. Nothing is predicted and nothing is decided.

  \item \textbf{Industrial validation} exists
  (\cite{holmegaard2016industrial,martins2018imdeld}, Theme~B), but the
  task is disaggregation, the methods are supervised, and the output is
  a retrospective attribution of energy rather than an operating state
  usable for control.

  \item \textbf{Quantification of idle-energy waste} exists (Theme~C),
  but rests on controlled laboratory measurement of individual
  processes and yields design guidance, not an online inference
  procedure applicable to a running facility.

  \item \textbf{Optimal shutdown decisions} exist (Theme~D), and the
  break-even formulation is well understood, but every such policy
  takes the idle duration as a \emph{given} input---supplied by a
  deterministic schedule or an analyst-calibrated arrival
  process---and is evaluated in simulation.
\end{itemize}

\medskip
\noindent\fbox{\parbox{0.97\columnwidth}{%
\textbf{Research gap.} No existing work combines
(i)~label-free, temporally coherent inference of industrial machine
states from electrical measurements alone, (ii)~forecasting of the
\emph{duration} of the resulting non-productive intervals, and
(iii)~an economically optimised shutdown decision derived from that
forecast, validated end-to-end on real industrial measurement data.
The state-inference and decision-making capabilities have matured
independently and remain unconnected: the former produces labels no
controller consumes, the latter consumes durations no inference
procedure supplies.}}
\medskip

\noindent
Closing this gap requires more than concatenating existing components.
Three problems arise specifically at the interfaces. First, because no
ground-truth state annotation exists for industrial data, the inferred
partition cannot be validated by conventional supervised metrics, and an
alternative validation basis---independent physical and operational
evidence---must be constructed. Second, an i.i.d.\ clustering produces a
label sequence whose flicker makes duration forecasting meaningless,
so temporal coherence is a prerequisite for the forecasting stage rather
than a refinement of the inference stage. Third, a forecast duration
carries uncertainty that a deterministic break-even threshold ignores;
the decision layer must therefore be posed as a constrained
cost-minimisation problem in which the break-even point emerges as a
derived quantity rather than being supplied as a tuning parameter.
Section~\ref{sec:method} addresses each of these in turn.

%  Requires in the preamble:
%      \usepackage{booktabs}
%      \usepackage{threeparttable}   % for the table notes
%      \usepackage{multirow}
%      \usepackage{wasysym}          % \CIRCLE \LEFTcircle \Circle
%  Note: wasysym redefines \Box/\diamond; if that clashes with your
%  template, load it as \usepackage[nointegrals]{wasysym} or swap the
%  harvey-ball glyphs for {\ding{108}}/{\ding{109}} from pifont.
%  Bibliography: paper/references.bib
% =====================================================================

\section{Related Work}
\label{sec:related}

The problem addressed in this paper sits at the intersection of four
research communities that have so far developed largely in isolation:
(i)~unsupervised identification of machine operating states from
electrical measurements, (ii)~non-intrusive load monitoring (NILM),
(iii)~industrial energy optimisation at the process and system level,
and (iv)~predictive shutdown and decision support for manufacturing
equipment. We review each in turn, then summarise the resulting gap.


% ---------------------------------------------------------------------
\subsection{Machine State Identification}
\label{sec:rw:states}

The recognition that an electrically driven machine occupies a small
number of discrete operating regimes---typically \textsc{off},
\textsc{standby} (energised but not producing), and \textsc{working}---is
the foundation of essentially all energy-state analysis. ISO~14955-1
\cite{iso14955} formalises this decomposition as a design-time
requirement for energy-efficient machine tools, defining standby as the
operating mode in which a machine is ready for production but performs
no value-adding function. This definition is the one adopted throughout
the present work.

\subsubsection{Supervised approaches}
Where per-timestamp annotations are available, state identification
reduces to a standard classification problem, and support-vector
machines, random forests, and convolutional classifiers all perform
well. The practical obstacle in industrial deployments is that such
annotations essentially never exist: obtaining them requires either
instrumenting the machine's contactor and PLC directly or having an
operator observe each machine over an extended period. Both are
prohibitively expensive at facility scale, and neither transfers to a
new machine without repeating the effort. Supervised methods are
therefore poorly matched to the retrofit scenario that motivates this
paper.

\subsubsection{Unsupervised approaches}
Unsupervised state inference avoids the annotation cost entirely.
$k$-means \cite{lloyd1982least} is the simplest option but imposes
isotropic, equal-variance clusters, which is a poor model for
electrical states whose dispersion grows with load. Gaussian mixture
models (GMM) \cite{mclachlan2000finite}, fitted by
expectation--maximisation \cite{dempster1977em}, relax both assumptions:
each state receives its own mean and full covariance, and assignments
are soft, which matters precisely at the \textsc{standby}/\textsc{working}
boundary where the two regimes overlap. Density-based methods such as
DBSCAN \cite{ester1996dbscan} and HDBSCAN \cite{campello2015hdbscan}
make no parametric shape assumption and identify noise points
explicitly, while spectral clustering \cite{ng2001spectral} operates on
the affinity graph and can recover non-convex structure. Model order is
conventionally selected by the Bayesian information criterion
\cite{schwarz1978bic} or by silhouette analysis
\cite{rousseeuw1987silhouettes}.

Neither criterion selects a state count on records of the size
considered here, and this is worth stating because it determines the
design of the method proposed below. The BIC penalty grows as
$\tfrac{1}{2}p\log n$ while the likelihood gained by splitting a
non-Gaussian component grows linearly in $n$, so at $n>10^{6}$ samples
per machine the penalty is negligible and a mixture model keeps buying
components in order to approximate a distribution that is not a mixture
of $k$ Gaussians for any $k$. Scanning $k\in[2,8]$ over the eight
machines of Section~\ref{sec:multimachine}, BIC is still falling at the
boundary on seven of them, and the improvement from $k=4$ to $k=8$ is
$3.1$--$33.1\%$ of the criterion's own value. Silhouette behaves
similarly. Statistical order selection must therefore be supplemented
by an external, physical admissibility criterion rather than merely
confirmed by one --- which is the role of the validation battery of
Section~\ref{sec:method:accept_reject}, and the reason it is part of
the method here rather than an appendix to it.

\subsubsection{Temporal structure}
A limitation common to all of the above is that they treat samples as
independent and identically distributed. Applied to a 1\,Hz electrical
stream, an i.i.d.\ clustering assigns each second in isolation and
consequently produces rapid, physically impossible state alternation
around cluster boundaries---a failure mode we quantify as the
\emph{flicker rate} in Section~\ref{sec:results}. Hidden Markov models
\cite{rabiner1989tutorial} address this by making the state sequence
itself the object of inference: the transition matrix encodes state
persistence, and Viterbi decoding returns the globally most probable
sequence rather than a concatenation of independent decisions. Fox
\textit{et al.} \cite{fox2011sticky} show that an explicit bias toward
self-transition (a ``sticky'' prior) reduces over-segmentation when the
emission distributions overlap, which is the regime encountered here.

It is worth stating precisely what this machinery can and cannot do,
because the point is routinely overstated. In a Viterbi decode the cost
of changing state is $\log(p_{\text{self}}/p_{\text{switch}})$, which is
a few nats even for extremely persistent transition matrices and is
bounded above however strong the prior is made. The emission term
carries no such bound: full-covariance Gaussians fitted to low-noise,
multi-channel electrical measurements can differ by hundreds of nats
between adjacent states. Where that holds, the transition model is
arithmetically unable to influence the decoded path and the HMM
degenerates toward per-sample assignment. Guaranteeing a minimum state
duration therefore requires modelling duration explicitly---a hidden
semi-Markov formulation \cite{yu2010hsmm}, or an equivalent
duration constraint imposed at decode time---rather than a heavier
transition prior. Section~\ref{sec:results} quantifies this for the
present data.

\subsubsection{Threshold-free decision rules}
Where a scalar decision boundary is unavoidable---for instance when
converting a continuous power-factor distribution into a binary
load/no-load verdict---Otsu's criterion \cite{otsu1979threshold}
selects the threshold that maximises between-class variance, removing
the arbitrariness of a hand-tuned cut-off. Agreement among several
independently derived partitions can then be aggregated using cluster
ensemble methods \cite{strehl2002cluster}.


% ---------------------------------------------------------------------
\subsection{Non-Intrusive Load Monitoring}
\label{sec:rw:nilm}

NILM originates with Hart \cite{hart1992nonintrusive}, who showed that
individual appliance activity can be recovered from an aggregate meter
by detecting steady-state step changes in real and reactive power.
Subsequent surveys \cite{zeifman2011nilm,schirmer2023nilmreview} trace
the field's progression from event-based signature matching to
probabilistic and, more recently, deep-learning formulations. Kolter
and Jaakkola \cite{kolter2012afamap} cast disaggregation as inference in
an additive factorial HMM, making explicit the connection between NILM
and the temporal state models of Section~\ref{sec:rw:states}. The deep
learning era is represented by Neural NILM \cite{kelly2015neuralnilm},
sequence-to-point regression \cite{zhang2018seq2point}, and denoising
autoencoders \cite{bonfigli2018denoising}. Standardised toolkits
\cite{batra2014nilmtk} and evaluation protocols
\cite{makonin2015performance} have made results across this line of
work broadly comparable.

Two properties of this literature are decisive for the present paper.
First, it is overwhelmingly \emph{residential}. The canonical benchmarks
---REDD \cite{kolter2011redd} and UK-DALE---consist of household
appliances whose signatures are sparse, short, and well separated.
Industrial machinery behaves very differently: loads are large,
continuously modulated by process throughput, dominated by induction
motors whose reactive power draw persists at zero mechanical output, and
governed by shift schedules rather than occupant behaviour. Holmegaard
and Kj\ae{}rgaard \cite{holmegaard2016industrial} evaluate established
NILM algorithms on industrial loads and report substantially degraded
performance, attributing it to precisely these differences. Martins
\textit{et al.} \cite{martins2018imdeld} apply a deep generative model
to industrial disaggregation and, in doing so, release the IMDELD
dataset \cite{imdeld2018dataset}---eleven meters recording eight heavy
machines in a Brazilian poultry-feed pellet plant at 1\,Hz across a
five-month window, of which the covered fraction differs by machine
(Section~\ref{sec:method:data}). IMDELD is the benchmark used throughout this
paper; its authors characterise the machines as three-state systems but
supply no per-timestamp state annotation, which is the reason the
present work must proceed without labels and validate its partition
against physics instead.

Second, and more fundamentally, the objective of NILM is
\emph{attribution}: given an aggregate signal, determine which appliance
consumed what. It is a retrospective accounting exercise. NILM does not
forecast, and it does not act. No NILM system, industrial or
residential, closes the loop from an inferred state to an actuation
decision. That step belongs to the two literatures reviewed next.


% ---------------------------------------------------------------------
\subsection{Industrial Energy Optimisation}
\label{sec:rw:industrial}

The manufacturing engineering community has independently established
that non-productive energy dominates the energy budget of discrete-part
production. Gutowski \textit{et al.} \cite{gutowski2006electrical} and
Dahmus and Gutowski \cite{dahmus2004environmental} show that the
load-independent baseline of a machine tool---spindle idling, hydraulics,
coolant pumps, control electronics---frequently exceeds the energy
actually delivered to the workpiece, and that this fraction rises as
automation increases. Kara and Li \cite{kara2011unit} formalise the
relationship as a unit-process model in which specific energy
consumption is inversely related to material removal rate, so that a
lightly loaded or idle machine is intrinsically inefficient. Devoldere
\textit{et al.} \cite{devoldere2007improvement} quantify the resulting
improvement potential and identify idle-mode elimination as the single
largest lever available without redesigning the machine. Duflou
\textit{et al.} \cite{duflou2012towards} synthesise this body of work
into a systems-level roadmap, placing operational measures---switching
equipment off when it is not needed---at the machine and cell levels of
their hierarchy.

On the management side, ISO~50001 \cite{iso50001} provides the
plan--do--check--act framework within which energy performance
indicators are defined and tracked, and ISO~14955-1 \cite{iso14955}
supplies the operating-mode taxonomy that makes standby energy a
measurable quantity in the first place.

The limitation of this literature is methodological rather than
conceptual. Its energy models are typically built from controlled
laboratory measurements of a single machine under prescribed process
parameters, and its recommendations are design guidelines and audit
procedures. It does not address how a facility, given only the
electrical measurements it already collects and no process-parameter
instrumentation, should determine \emph{which} of its machines are idle,
\emph{when}, and for \emph{how long}. That inference problem is left
open.


% ---------------------------------------------------------------------
\subsection{Predictive Shutdown and Decision Support}
\label{sec:rw:decision}

The operations research community has studied the shutdown decision
itself in considerable depth. Mouzon \textit{et al.}
\cite{mouzon2007operational} establish the central trade-off: turning a
machine off saves its idle power but incurs a restart energy penalty and
a restart delay, so shutdown is economical only when the idle period
exceeds a break-even duration determined by the ratio of restart energy
to idle power. They propose dispatching rules and a mathematical
programming formulation that exploit this, and extend the treatment to
joint energy--tardiness objectives \cite{mouzon2008framework}. Chen
\textit{et al.} \cite{chen2013schedule} develop schedule-based operation
policies for production lines, Sun and Li \cite{sun2013opportunity}
estimate real-time control opportunities from buffer state, Li and Sun
\cite{li2016energy} formulate dynamic energy control as a Markov
decision process, Frigerio and Matta \cite{frigerio2015energy} derive
control strategies for machine tools under stochastic job arrivals, and
He \textit{et al.} \cite{he2015energy} embed energy responsiveness in
job-shop machine selection and sequencing.

This literature is mathematically mature, and the break-even
formulation of Mouzon \textit{et al.} is the natural objective for the
decision layer of any standby-elimination system. Its universal
assumption, however, is that the idle interval is \emph{known}---either
because it is dictated by a deterministic production schedule, or
because it is drawn from an analytically tractable arrival process
calibrated by the modeller. The policies are evaluated in simulation or
on analytical queueing models. What is absent is any mechanism for
obtaining the idle-duration input from a machine's own measured
electrical behaviour, which is the only signal available in a retrofit
installation. In other words: this community can decide optimally given
a duration, and the NILM community can infer a state from
measurements, but no work connects the two.


% ---------------------------------------------------------------------
\subsection{Synthesis and Research Gap}
\label{sec:rw:gap}

Table~\ref{tab:related} compares representative works from the four
themes against the capabilities required for an autonomous
standby-elimination system: whether the method infers machine state
without labels, whether it models temporal structure explicitly, whether
it forecasts the \emph{duration} of a future non-productive interval,
whether it converts that forecast into an economically justified control
decision, and whether it is validated on real industrial measurements.

\begin{table*}[!t]
\centering
\caption{Comparison of representative prior work against the
capabilities required for autonomous standby elimination.
\CIRCLE~= provided, \LEFTcircle~= partially provided, \Circle~= not
provided, --~= not applicable.}
\label{tab:related}
\begin{threeparttable}
\small
\setlength{\tabcolsep}{4pt}
\begin{tabular}{@{}lllccccc@{}}
\toprule
\multirow{2}{*}{\textbf{Work}} &
\multirow{2}{*}{\textbf{Method family}} &
\multirow{2}{*}{\textbf{Data domain}} &
\textbf{Labels} &
\textbf{Temporal} &
\textbf{Duration} &
\textbf{Economic} &
\textbf{Real ind.} \\
 & & & \textbf{required} & \textbf{model} & \textbf{forecast} &
 \textbf{decision} & \textbf{data} \\
\midrule
\multicolumn{8}{@{}l}{\textit{Theme A --- machine state identification}}\\
Lloyd \cite{lloyd1982least}          & $k$-means                    & generic                & no  & \Circle    & \Circle & \Circle & \Circle \\
McLachlan \& Peel \cite{mclachlan2000finite} & Gaussian mixture     & generic                & no  & \Circle    & \Circle & \Circle & \Circle \\
Ester \textit{et al.} \cite{ester1996dbscan} & DBSCAN               & generic                & no  & \Circle    & \Circle & \Circle & \Circle \\
Fox \textit{et al.} \cite{fox2011sticky}     & sticky HDP-HMM       & audio                  & no  & \CIRCLE    & \Circle & \Circle & \Circle \\
\addlinespace
\multicolumn{8}{@{}l}{\textit{Theme B --- non-intrusive load monitoring}}\\
Hart \cite{hart1992nonintrusive}     & event/steady-state signature & residential            & yes\tnote{a} & \LEFTcircle & \Circle & \Circle & \Circle \\
Kolter \& Jaakkola \cite{kolter2012afamap} & additive factorial HMM & REDD (residential)     & no  & \CIRCLE    & \Circle & \Circle & \Circle \\
Kelly \& Knottenbelt \cite{kelly2015neuralnilm} & deep NN (dAE, LSTM) & UK-DALE (residential) & yes & \LEFTcircle & \Circle & \Circle & \Circle \\
Zhang \textit{et al.} \cite{zhang2018seq2point} & seq2point CNN     & UK-DALE, REDD          & yes & \LEFTcircle & \Circle & \Circle & \Circle \\
Holmegaard \& Kj\ae{}rgaard \cite{holmegaard2016industrial} & NILM algorithm evaluation & industrial & yes & \LEFTcircle & \Circle & \Circle & \CIRCLE \\
Martins \textit{et al.} \cite{martins2018imdeld} & deep generative disaggregation & IMDELD (industrial) & yes & \LEFTcircle & \Circle & \Circle & \CIRCLE \\
\addlinespace
\multicolumn{8}{@{}l}{\textit{Theme C --- industrial energy optimisation}}\\
Gutowski \textit{et al.} \cite{gutowski2006electrical} & thermodynamic/empirical process model & lab machine tools & -- & \Circle & \Circle & \Circle & \LEFTcircle \\
Kara \& Li \cite{kara2011unit}       & unit-process energy model    & lab machine tools      & --  & \Circle    & \Circle & \Circle & \LEFTcircle \\
Duflou \textit{et al.} \cite{duflou2012towards} & systems-level roadmap & survey             & --  & \Circle    & \Circle & \LEFTcircle & \LEFTcircle \\
\addlinespace
\multicolumn{8}{@{}l}{\textit{Theme D --- predictive shutdown and decision support}}\\
Mouzon \textit{et al.} \cite{mouzon2007operational} & dispatching rules, MILP & simulated single machine & -- & \Circle & \Circle & \CIRCLE & \Circle \\
Chen \textit{et al.} \cite{chen2013schedule} & schedule-based control  & simulated line     & --  & \LEFTcircle & \Circle & \CIRCLE & \Circle \\
Frigerio \& Matta \cite{frigerio2015energy} & stochastic control policy & analytical model  & --  & \CIRCLE    & \LEFTcircle\tnote{b} & \CIRCLE & \Circle \\
He \textit{et al.} \cite{he2015energy} & energy-responsive scheduling & job-shop model      & --  & \Circle    & \Circle & \CIRCLE & \Circle \\
\midrule
\textbf{This work}                   & \textbf{physics-gated GMM--HMM $\rightarrow$ duration forecast $\rightarrow$ optimisation} & \textbf{IMDELD (industrial)} & \textbf{no} & \CIRCLE & \CIRCLE & \CIRCLE & \CIRCLE \\
\bottomrule
\end{tabular}
\begin{tablenotes}[flushleft]\footnotesize
\item[a] Hart's formulation is unsupervised in its clustering step but
requires a manually curated appliance signature library to name the
recovered clusters.
\item[b] Idle duration is modelled as a known stochastic arrival
process calibrated by the analyst, not forecast from measured machine
behaviour.
\end{tablenotes}
\end{threeparttable}
\end{table*}

Reading Table~\ref{tab:related} column-wise makes the gap explicit.
Every capability required for autonomous standby elimination has been
solved somewhere, but the last column of any single row is always
empty:

\begin{itemize}
  \item \textbf{Label-free state inference with temporal structure}
  exists (Theme~A, and \cite{kolter2012afamap} in Theme~B), but is
  demonstrated on generic or residential data and terminates at the
  labelled sequence. Nothing is predicted and nothing is decided.

  \item \textbf{Industrial validation} exists
  (\cite{holmegaard2016industrial,martins2018imdeld}, Theme~B), but the
  task is disaggregation, the methods are supervised, and the output is
  a retrospective attribution of energy rather than an operating state
  usable for control.

  \item \textbf{Quantification of idle-energy waste} exists (Theme~C),
  but rests on controlled laboratory measurement of individual
  processes and yields design guidance, not an online inference
  procedure applicable to a running facility.

  \item \textbf{Optimal shutdown decisions} exist (Theme~D), and the
  break-even formulation is well understood, but every such policy
  takes the idle duration as a \emph{given} input---supplied by a
  deterministic schedule or an analyst-calibrated arrival
  process---and is evaluated in simulation.
\end{itemize}

\medskip
\noindent\fbox{\parbox{0.97\columnwidth}{%
\textbf{Research gap.} No existing work combines
(i)~label-free, temporally coherent inference of industrial machine
states from electrical measurements alone, (ii)~forecasting of the
\emph{duration} of the resulting non-productive intervals, and
(iii)~an economically optimised shutdown decision derived from that
forecast, validated end-to-end on real industrial measurement data.
The state-inference and decision-making capabilities have matured
independently and remain unconnected: the former produces labels no
controller consumes, the latter consumes durations no inference
procedure supplies.}}
\medskip

\noindent
Closing this gap requires more than concatenating existing components.
Three problems arise specifically at the interfaces. First, because no
ground-truth state annotation exists for industrial data, the inferred
partition cannot be validated by conventional supervised metrics, and an
alternative validation basis---independent physical and operational
evidence---must be constructed. Second, an i.i.d.\ clustering produces a
label sequence whose flicker makes duration forecasting meaningless,
so temporal coherence is a prerequisite for the forecasting stage rather
than a refinement of the inference stage. Third, a forecast duration
carries uncertainty that a deterministic break-even threshold ignores;
the decision layer must therefore be posed as a constrained
cost-minimisation problem in which the break-even point emerges as a
derived quantity rather than being supplied as a tuning parameter.
Section~\ref{sec:method} addresses each of these in turn.

%  Requires in the preamble:
%      \usepackage{booktabs}
%      \usepackage{threeparttable}   % for the table notes
%      \usepackage{multirow}
%      \usepackage{wasysym}          % \CIRCLE \LEFTcircle \Circle
%  Note: wasysym redefines \Box/\diamond; if that clashes with your
%  template, load it as \usepackage[nointegrals]{wasysym} or swap the
%  harvey-ball glyphs for {\ding{108}}/{\ding{109}} from pifont.
%  Bibliography: paper/references.bib
% =====================================================================

\section{Related Work}
\label{sec:related}

The problem addressed in this paper sits at the intersection of four
research communities that have so far developed largely in isolation:
(i)~unsupervised identification of machine operating states from
electrical measurements, (ii)~non-intrusive load monitoring (NILM),
(iii)~industrial energy optimisation at the process and system level,
and (iv)~predictive shutdown and decision support for manufacturing
equipment. We review each in turn, then summarise the resulting gap.


% ---------------------------------------------------------------------
\subsection{Machine State Identification}
\label{sec:rw:states}

The recognition that an electrically driven machine occupies a small
number of discrete operating regimes---typically \textsc{off},
\textsc{standby} (energised but not producing), and \textsc{working}---is
the foundation of essentially all energy-state analysis. ISO~14955-1
\cite{iso14955} formalises this decomposition as a design-time
requirement for energy-efficient machine tools, defining standby as the
operating mode in which a machine is ready for production but performs
no value-adding function. This definition is the one adopted throughout
the present work.

\subsubsection{Supervised approaches}
Where per-timestamp annotations are available, state identification
reduces to a standard classification problem, and support-vector
machines, random forests, and convolutional classifiers all perform
well. The practical obstacle in industrial deployments is that such
annotations essentially never exist: obtaining them requires either
instrumenting the machine's contactor and PLC directly or having an
operator observe each machine over an extended period. Both are
prohibitively expensive at facility scale, and neither transfers to a
new machine without repeating the effort. Supervised methods are
therefore poorly matched to the retrofit scenario that motivates this
paper.

\subsubsection{Unsupervised approaches}
Unsupervised state inference avoids the annotation cost entirely.
$k$-means \cite{lloyd1982least} is the simplest option but imposes
isotropic, equal-variance clusters, which is a poor model for
electrical states whose dispersion grows with load. Gaussian mixture
models (GMM) \cite{mclachlan2000finite}, fitted by
expectation--maximisation \cite{dempster1977em}, relax both assumptions:
each state receives its own mean and full covariance, and assignments
are soft, which matters precisely at the \textsc{standby}/\textsc{working}
boundary where the two regimes overlap. Density-based methods such as
DBSCAN \cite{ester1996dbscan} and HDBSCAN \cite{campello2015hdbscan}
make no parametric shape assumption and identify noise points
explicitly, while spectral clustering \cite{ng2001spectral} operates on
the affinity graph and can recover non-convex structure. Model order is
conventionally selected by the Bayesian information criterion
\cite{schwarz1978bic} or by silhouette analysis
\cite{rousseeuw1987silhouettes}.

Neither criterion selects a state count on records of the size
considered here, and this is worth stating because it determines the
design of the method proposed below. The BIC penalty grows as
$\tfrac{1}{2}p\log n$ while the likelihood gained by splitting a
non-Gaussian component grows linearly in $n$, so at $n>10^{6}$ samples
per machine the penalty is negligible and a mixture model keeps buying
components in order to approximate a distribution that is not a mixture
of $k$ Gaussians for any $k$. Scanning $k\in[2,8]$ over the eight
machines of Section~\ref{sec:multimachine}, BIC is still falling at the
boundary on seven of them, and the improvement from $k=4$ to $k=8$ is
$3.1$--$33.1\%$ of the criterion's own value. Silhouette behaves
similarly. Statistical order selection must therefore be supplemented
by an external, physical admissibility criterion rather than merely
confirmed by one --- which is the role of the validation battery of
Section~\ref{sec:method:accept_reject}, and the reason it is part of
the method here rather than an appendix to it.

\subsubsection{Temporal structure}
A limitation common to all of the above is that they treat samples as
independent and identically distributed. Applied to a 1\,Hz electrical
stream, an i.i.d.\ clustering assigns each second in isolation and
consequently produces rapid, physically impossible state alternation
around cluster boundaries---a failure mode we quantify as the
\emph{flicker rate} in Section~\ref{sec:results}. Hidden Markov models
\cite{rabiner1989tutorial} address this by making the state sequence
itself the object of inference: the transition matrix encodes state
persistence, and Viterbi decoding returns the globally most probable
sequence rather than a concatenation of independent decisions. Fox
\textit{et al.} \cite{fox2011sticky} show that an explicit bias toward
self-transition (a ``sticky'' prior) reduces over-segmentation when the
emission distributions overlap, which is the regime encountered here.

It is worth stating precisely what this machinery can and cannot do,
because the point is routinely overstated. In a Viterbi decode the cost
of changing state is $\log(p_{\text{self}}/p_{\text{switch}})$, which is
a few nats even for extremely persistent transition matrices and is
bounded above however strong the prior is made. The emission term
carries no such bound: full-covariance Gaussians fitted to low-noise,
multi-channel electrical measurements can differ by hundreds of nats
between adjacent states. Where that holds, the transition model is
arithmetically unable to influence the decoded path and the HMM
degenerates toward per-sample assignment. Guaranteeing a minimum state
duration therefore requires modelling duration explicitly---a hidden
semi-Markov formulation \cite{yu2010hsmm}, or an equivalent
duration constraint imposed at decode time---rather than a heavier
transition prior. Section~\ref{sec:results} quantifies this for the
present data.

\subsubsection{Threshold-free decision rules}
Where a scalar decision boundary is unavoidable---for instance when
converting a continuous power-factor distribution into a binary
load/no-load verdict---Otsu's criterion \cite{otsu1979threshold}
selects the threshold that maximises between-class variance, removing
the arbitrariness of a hand-tuned cut-off. Agreement among several
independently derived partitions can then be aggregated using cluster
ensemble methods \cite{strehl2002cluster}.


% ---------------------------------------------------------------------
\subsection{Non-Intrusive Load Monitoring}
\label{sec:rw:nilm}

NILM originates with Hart \cite{hart1992nonintrusive}, who showed that
individual appliance activity can be recovered from an aggregate meter
by detecting steady-state step changes in real and reactive power.
Subsequent surveys \cite{zeifman2011nilm,schirmer2023nilmreview} trace
the field's progression from event-based signature matching to
probabilistic and, more recently, deep-learning formulations. Kolter
and Jaakkola \cite{kolter2012afamap} cast disaggregation as inference in
an additive factorial HMM, making explicit the connection between NILM
and the temporal state models of Section~\ref{sec:rw:states}. The deep
learning era is represented by Neural NILM \cite{kelly2015neuralnilm},
sequence-to-point regression \cite{zhang2018seq2point}, and denoising
autoencoders \cite{bonfigli2018denoising}. Standardised toolkits
\cite{batra2014nilmtk} and evaluation protocols
\cite{makonin2015performance} have made results across this line of
work broadly comparable.

Two properties of this literature are decisive for the present paper.
First, it is overwhelmingly \emph{residential}. The canonical benchmarks
---REDD \cite{kolter2011redd} and UK-DALE---consist of household
appliances whose signatures are sparse, short, and well separated.
Industrial machinery behaves very differently: loads are large,
continuously modulated by process throughput, dominated by induction
motors whose reactive power draw persists at zero mechanical output, and
governed by shift schedules rather than occupant behaviour. Holmegaard
and Kj\ae{}rgaard \cite{holmegaard2016industrial} evaluate established
NILM algorithms on industrial loads and report substantially degraded
performance, attributing it to precisely these differences. Martins
\textit{et al.} \cite{martins2018imdeld} apply a deep generative model
to industrial disaggregation and, in doing so, release the IMDELD
dataset \cite{imdeld2018dataset}---eleven meters recording eight heavy
machines in a Brazilian poultry-feed pellet plant at 1\,Hz across a
five-month window, of which the covered fraction differs by machine
(Section~\ref{sec:method:data}). IMDELD is the benchmark used throughout this
paper; its authors characterise the machines as three-state systems but
supply no per-timestamp state annotation, which is the reason the
present work must proceed without labels and validate its partition
against physics instead.

Second, and more fundamentally, the objective of NILM is
\emph{attribution}: given an aggregate signal, determine which appliance
consumed what. It is a retrospective accounting exercise. NILM does not
forecast, and it does not act. No NILM system, industrial or
residential, closes the loop from an inferred state to an actuation
decision. That step belongs to the two literatures reviewed next.


% ---------------------------------------------------------------------
\subsection{Industrial Energy Optimisation}
\label{sec:rw:industrial}

The manufacturing engineering community has independently established
that non-productive energy dominates the energy budget of discrete-part
production. Gutowski \textit{et al.} \cite{gutowski2006electrical} and
Dahmus and Gutowski \cite{dahmus2004environmental} show that the
load-independent baseline of a machine tool---spindle idling, hydraulics,
coolant pumps, control electronics---frequently exceeds the energy
actually delivered to the workpiece, and that this fraction rises as
automation increases. Kara and Li \cite{kara2011unit} formalise the
relationship as a unit-process model in which specific energy
consumption is inversely related to material removal rate, so that a
lightly loaded or idle machine is intrinsically inefficient. Devoldere
\textit{et al.} \cite{devoldere2007improvement} quantify the resulting
improvement potential and identify idle-mode elimination as the single
largest lever available without redesigning the machine. Duflou
\textit{et al.} \cite{duflou2012towards} synthesise this body of work
into a systems-level roadmap, placing operational measures---switching
equipment off when it is not needed---at the machine and cell levels of
their hierarchy.

On the management side, ISO~50001 \cite{iso50001} provides the
plan--do--check--act framework within which energy performance
indicators are defined and tracked, and ISO~14955-1 \cite{iso14955}
supplies the operating-mode taxonomy that makes standby energy a
measurable quantity in the first place.

The limitation of this literature is methodological rather than
conceptual. Its energy models are typically built from controlled
laboratory measurements of a single machine under prescribed process
parameters, and its recommendations are design guidelines and audit
procedures. It does not address how a facility, given only the
electrical measurements it already collects and no process-parameter
instrumentation, should determine \emph{which} of its machines are idle,
\emph{when}, and for \emph{how long}. That inference problem is left
open.


% ---------------------------------------------------------------------
\subsection{Predictive Shutdown and Decision Support}
\label{sec:rw:decision}

The operations research community has studied the shutdown decision
itself in considerable depth. Mouzon \textit{et al.}
\cite{mouzon2007operational} establish the central trade-off: turning a
machine off saves its idle power but incurs a restart energy penalty and
a restart delay, so shutdown is economical only when the idle period
exceeds a break-even duration determined by the ratio of restart energy
to idle power. They propose dispatching rules and a mathematical
programming formulation that exploit this, and extend the treatment to
joint energy--tardiness objectives \cite{mouzon2008framework}. Chen
\textit{et al.} \cite{chen2013schedule} develop schedule-based operation
policies for production lines, Sun and Li \cite{sun2013opportunity}
estimate real-time control opportunities from buffer state, Li and Sun
\cite{li2016energy} formulate dynamic energy control as a Markov
decision process, Frigerio and Matta \cite{frigerio2015energy} derive
control strategies for machine tools under stochastic job arrivals, and
He \textit{et al.} \cite{he2015energy} embed energy responsiveness in
job-shop machine selection and sequencing.

This literature is mathematically mature, and the break-even
formulation of Mouzon \textit{et al.} is the natural objective for the
decision layer of any standby-elimination system. Its universal
assumption, however, is that the idle interval is \emph{known}---either
because it is dictated by a deterministic production schedule, or
because it is drawn from an analytically tractable arrival process
calibrated by the modeller. The policies are evaluated in simulation or
on analytical queueing models. What is absent is any mechanism for
obtaining the idle-duration input from a machine's own measured
electrical behaviour, which is the only signal available in a retrofit
installation. In other words: this community can decide optimally given
a duration, and the NILM community can infer a state from
measurements, but no work connects the two.


% ---------------------------------------------------------------------
\subsection{Synthesis and Research Gap}
\label{sec:rw:gap}

Table~\ref{tab:related} compares representative works from the four
themes against the capabilities required for an autonomous
standby-elimination system: whether the method infers machine state
without labels, whether it models temporal structure explicitly, whether
it forecasts the \emph{duration} of a future non-productive interval,
whether it converts that forecast into an economically justified control
decision, and whether it is validated on real industrial measurements.

\begin{table*}[!t]
\centering
\caption{Comparison of representative prior work against the
capabilities required for autonomous standby elimination.
\CIRCLE~= provided, \LEFTcircle~= partially provided, \Circle~= not
provided, --~= not applicable.}
\label{tab:related}
\begin{threeparttable}
\small
\setlength{\tabcolsep}{4pt}
\begin{tabular}{@{}lllccccc@{}}
\toprule
\multirow{2}{*}{\textbf{Work}} &
\multirow{2}{*}{\textbf{Method family}} &
\multirow{2}{*}{\textbf{Data domain}} &
\textbf{Labels} &
\textbf{Temporal} &
\textbf{Duration} &
\textbf{Economic} &
\textbf{Real ind.} \\
 & & & \textbf{required} & \textbf{model} & \textbf{forecast} &
 \textbf{decision} & \textbf{data} \\
\midrule
\multicolumn{8}{@{}l}{\textit{Theme A --- machine state identification}}\\
Lloyd \cite{lloyd1982least}          & $k$-means                    & generic                & no  & \Circle    & \Circle & \Circle & \Circle \\
McLachlan \& Peel \cite{mclachlan2000finite} & Gaussian mixture     & generic                & no  & \Circle    & \Circle & \Circle & \Circle \\
Ester \textit{et al.} \cite{ester1996dbscan} & DBSCAN               & generic                & no  & \Circle    & \Circle & \Circle & \Circle \\
Fox \textit{et al.} \cite{fox2011sticky}     & sticky HDP-HMM       & audio                  & no  & \CIRCLE    & \Circle & \Circle & \Circle \\
\addlinespace
\multicolumn{8}{@{}l}{\textit{Theme B --- non-intrusive load monitoring}}\\
Hart \cite{hart1992nonintrusive}     & event/steady-state signature & residential            & yes\tnote{a} & \LEFTcircle & \Circle & \Circle & \Circle \\
Kolter \& Jaakkola \cite{kolter2012afamap} & additive factorial HMM & REDD (residential)     & no  & \CIRCLE    & \Circle & \Circle & \Circle \\
Kelly \& Knottenbelt \cite{kelly2015neuralnilm} & deep NN (dAE, LSTM) & UK-DALE (residential) & yes & \LEFTcircle & \Circle & \Circle & \Circle \\
Zhang \textit{et al.} \cite{zhang2018seq2point} & seq2point CNN     & UK-DALE, REDD          & yes & \LEFTcircle & \Circle & \Circle & \Circle \\
Holmegaard \& Kj\ae{}rgaard \cite{holmegaard2016industrial} & NILM algorithm evaluation & industrial & yes & \LEFTcircle & \Circle & \Circle & \CIRCLE \\
Martins \textit{et al.} \cite{martins2018imdeld} & deep generative disaggregation & IMDELD (industrial) & yes & \LEFTcircle & \Circle & \Circle & \CIRCLE \\
\addlinespace
\multicolumn{8}{@{}l}{\textit{Theme C --- industrial energy optimisation}}\\
Gutowski \textit{et al.} \cite{gutowski2006electrical} & thermodynamic/empirical process model & lab machine tools & -- & \Circle & \Circle & \Circle & \LEFTcircle \\
Kara \& Li \cite{kara2011unit}       & unit-process energy model    & lab machine tools      & --  & \Circle    & \Circle & \Circle & \LEFTcircle \\
Duflou \textit{et al.} \cite{duflou2012towards} & systems-level roadmap & survey             & --  & \Circle    & \Circle & \LEFTcircle & \LEFTcircle \\
\addlinespace
\multicolumn{8}{@{}l}{\textit{Theme D --- predictive shutdown and decision support}}\\
Mouzon \textit{et al.} \cite{mouzon2007operational} & dispatching rules, MILP & simulated single machine & -- & \Circle & \Circle & \CIRCLE & \Circle \\
Chen \textit{et al.} \cite{chen2013schedule} & schedule-based control  & simulated line     & --  & \LEFTcircle & \Circle & \CIRCLE & \Circle \\
Frigerio \& Matta \cite{frigerio2015energy} & stochastic control policy & analytical model  & --  & \CIRCLE    & \LEFTcircle\tnote{b} & \CIRCLE & \Circle \\
He \textit{et al.} \cite{he2015energy} & energy-responsive scheduling & job-shop model      & --  & \Circle    & \Circle & \CIRCLE & \Circle \\
\midrule
\textbf{This work}                   & \textbf{physics-gated GMM--HMM $\rightarrow$ duration forecast $\rightarrow$ optimisation} & \textbf{IMDELD (industrial)} & \textbf{no} & \CIRCLE & \CIRCLE & \CIRCLE & \CIRCLE \\
\bottomrule
\end{tabular}
\begin{tablenotes}[flushleft]\footnotesize
\item[a] Hart's formulation is unsupervised in its clustering step but
requires a manually curated appliance signature library to name the
recovered clusters.
\item[b] Idle duration is modelled as a known stochastic arrival
process calibrated by the analyst, not forecast from measured machine
behaviour.
\end{tablenotes}
\end{threeparttable}
\end{table*}

Reading Table~\ref{tab:related} column-wise makes the gap explicit.
Every capability required for autonomous standby elimination has been
solved somewhere, but the last column of any single row is always
empty:

\begin{itemize}
  \item \textbf{Label-free state inference with temporal structure}
  exists (Theme~A, and \cite{kolter2012afamap} in Theme~B), but is
  demonstrated on generic or residential data and terminates at the
  labelled sequence. Nothing is predicted and nothing is decided.

  \item \textbf{Industrial validation} exists
  (\cite{holmegaard2016industrial,martins2018imdeld}, Theme~B), but the
  task is disaggregation, the methods are supervised, and the output is
  a retrospective attribution of energy rather than an operating state
  usable for control.

  \item \textbf{Quantification of idle-energy waste} exists (Theme~C),
  but rests on controlled laboratory measurement of individual
  processes and yields design guidance, not an online inference
  procedure applicable to a running facility.

  \item \textbf{Optimal shutdown decisions} exist (Theme~D), and the
  break-even formulation is well understood, but every such policy
  takes the idle duration as a \emph{given} input---supplied by a
  deterministic schedule or an analyst-calibrated arrival
  process---and is evaluated in simulation.
\end{itemize}

\medskip
\noindent\fbox{\parbox{0.97\columnwidth}{%
\textbf{Research gap.} No existing work combines
(i)~label-free, temporally coherent inference of industrial machine
states from electrical measurements alone, (ii)~forecasting of the
\emph{duration} of the resulting non-productive intervals, and
(iii)~an economically optimised shutdown decision derived from that
forecast, validated end-to-end on real industrial measurement data.
The state-inference and decision-making capabilities have matured
independently and remain unconnected: the former produces labels no
controller consumes, the latter consumes durations no inference
procedure supplies.}}
\medskip

\noindent
Closing this gap requires more than concatenating existing components.
Three problems arise specifically at the interfaces. First, because no
ground-truth state annotation exists for industrial data, the inferred
partition cannot be validated by conventional supervised metrics, and an
alternative validation basis---independent physical and operational
evidence---must be constructed. Second, an i.i.d.\ clustering produces a
label sequence whose flicker makes duration forecasting meaningless,
so temporal coherence is a prerequisite for the forecasting stage rather
than a refinement of the inference stage. Third, a forecast duration
carries uncertainty that a deterministic break-even threshold ignores;
the decision layer must therefore be posed as a constrained
cost-minimisation problem in which the break-even point emerges as a
derived quantity rather than being supplied as a tuning parameter.
Section~\ref{sec:method} addresses each of these in turn.

%  Requires in the preamble:
%      \usepackage{booktabs}
%      \usepackage{threeparttable}   % for the table notes
%      \usepackage{multirow}
%      \usepackage{wasysym}          % \CIRCLE \LEFTcircle \Circle
%  Note: wasysym redefines \Box/\diamond; if that clashes with your
%  template, load it as \usepackage[nointegrals]{wasysym} or swap the
%  harvey-ball glyphs for {\ding{108}}/{\ding{109}} from pifont.
%  Bibliography: paper/references.bib
% =====================================================================

\section{Related Work}
\label{sec:related}

The problem addressed in this paper sits at the intersection of four
research communities that have so far developed largely in isolation:
(i)~unsupervised identification of machine operating states from
electrical measurements, (ii)~non-intrusive load monitoring (NILM),
(iii)~industrial energy optimisation at the process and system level,
and (iv)~predictive shutdown and decision support for manufacturing
equipment. We review each in turn, then summarise the resulting gap.


% ---------------------------------------------------------------------
\subsection{Machine State Identification}
\label{sec:rw:states}

The recognition that an electrically driven machine occupies a small
number of discrete operating regimes---typically \textsc{off},
\textsc{standby} (energised but not producing), and \textsc{working}---is
the foundation of essentially all energy-state analysis. ISO~14955-1
\cite{iso14955} formalises this decomposition as a design-time
requirement for energy-efficient machine tools, defining standby as the
operating mode in which a machine is ready for production but performs
no value-adding function. This definition is the one adopted throughout
the present work.

\subsubsection{Supervised approaches}
Where per-timestamp annotations are available, state identification
reduces to a standard classification problem, and support-vector
machines, random forests, and convolutional classifiers all perform
well. The practical obstacle in industrial deployments is that such
annotations essentially never exist: obtaining them requires either
instrumenting the machine's contactor and PLC directly or having an
operator observe each machine over an extended period. Both are
prohibitively expensive at facility scale, and neither transfers to a
new machine without repeating the effort. Supervised methods are
therefore poorly matched to the retrofit scenario that motivates this
paper.

\subsubsection{Unsupervised approaches}
Unsupervised state inference avoids the annotation cost entirely.
$k$-means \cite{lloyd1982least} is the simplest option but imposes
isotropic, equal-variance clusters, which is a poor model for
electrical states whose dispersion grows with load. Gaussian mixture
models (GMM) \cite{mclachlan2000finite}, fitted by
expectation--maximisation \cite{dempster1977em}, relax both assumptions:
each state receives its own mean and full covariance, and assignments
are soft, which matters precisely at the \textsc{standby}/\textsc{working}
boundary where the two regimes overlap. Density-based methods such as
DBSCAN \cite{ester1996dbscan} and HDBSCAN \cite{campello2015hdbscan}
make no parametric shape assumption and identify noise points
explicitly, while spectral clustering \cite{ng2001spectral} operates on
the affinity graph and can recover non-convex structure. Model order is
conventionally selected by the Bayesian information criterion
\cite{schwarz1978bic} or by silhouette analysis
\cite{rousseeuw1987silhouettes}.

Neither criterion selects a state count on records of the size
considered here, and this is worth stating because it determines the
design of the method proposed below. The BIC penalty grows as
$\tfrac{1}{2}p\log n$ while the likelihood gained by splitting a
non-Gaussian component grows linearly in $n$, so at $n>10^{6}$ samples
per machine the penalty is negligible and a mixture model keeps buying
components in order to approximate a distribution that is not a mixture
of $k$ Gaussians for any $k$. Scanning $k\in[2,8]$ over the eight
machines of Section~\ref{sec:multimachine}, BIC is still falling at the
boundary on seven of them, and the improvement from $k=4$ to $k=8$ is
$3.1$--$33.1\%$ of the criterion's own value. Silhouette behaves
similarly. Statistical order selection must therefore be supplemented
by an external, physical admissibility criterion rather than merely
confirmed by one --- which is the role of the validation battery of
Section~\ref{sec:method:accept_reject}, and the reason it is part of
the method here rather than an appendix to it.

\subsubsection{Temporal structure}
A limitation common to all of the above is that they treat samples as
independent and identically distributed. Applied to a 1\,Hz electrical
stream, an i.i.d.\ clustering assigns each second in isolation and
consequently produces rapid, physically impossible state alternation
around cluster boundaries---a failure mode we quantify as the
\emph{flicker rate} in Section~\ref{sec:results}. Hidden Markov models
\cite{rabiner1989tutorial} address this by making the state sequence
itself the object of inference: the transition matrix encodes state
persistence, and Viterbi decoding returns the globally most probable
sequence rather than a concatenation of independent decisions. Fox
\textit{et al.} \cite{fox2011sticky} show that an explicit bias toward
self-transition (a ``sticky'' prior) reduces over-segmentation when the
emission distributions overlap, which is the regime encountered here.

It is worth stating precisely what this machinery can and cannot do,
because the point is routinely overstated. In a Viterbi decode the cost
of changing state is $\log(p_{\text{self}}/p_{\text{switch}})$, which is
a few nats even for extremely persistent transition matrices and is
bounded above however strong the prior is made. The emission term
carries no such bound: full-covariance Gaussians fitted to low-noise,
multi-channel electrical measurements can differ by hundreds of nats
between adjacent states. Where that holds, the transition model is
arithmetically unable to influence the decoded path and the HMM
degenerates toward per-sample assignment. Guaranteeing a minimum state
duration therefore requires modelling duration explicitly---a hidden
semi-Markov formulation \cite{yu2010hsmm}, or an equivalent
duration constraint imposed at decode time---rather than a heavier
transition prior. Section~\ref{sec:results} quantifies this for the
present data.

\subsubsection{Threshold-free decision rules}
Where a scalar decision boundary is unavoidable---for instance when
converting a continuous power-factor distribution into a binary
load/no-load verdict---Otsu's criterion \cite{otsu1979threshold}
selects the threshold that maximises between-class variance, removing
the arbitrariness of a hand-tuned cut-off. Agreement among several
independently derived partitions can then be aggregated using cluster
ensemble methods \cite{strehl2002cluster}.


% ---------------------------------------------------------------------
\subsection{Non-Intrusive Load Monitoring}
\label{sec:rw:nilm}

NILM originates with Hart \cite{hart1992nonintrusive}, who showed that
individual appliance activity can be recovered from an aggregate meter
by detecting steady-state step changes in real and reactive power.
Subsequent surveys \cite{zeifman2011nilm,schirmer2023nilmreview} trace
the field's progression from event-based signature matching to
probabilistic and, more recently, deep-learning formulations. Kolter
and Jaakkola \cite{kolter2012afamap} cast disaggregation as inference in
an additive factorial HMM, making explicit the connection between NILM
and the temporal state models of Section~\ref{sec:rw:states}. The deep
learning era is represented by Neural NILM \cite{kelly2015neuralnilm},
sequence-to-point regression \cite{zhang2018seq2point}, and denoising
autoencoders \cite{bonfigli2018denoising}. Standardised toolkits
\cite{batra2014nilmtk} and evaluation protocols
\cite{makonin2015performance} have made results across this line of
work broadly comparable.

Two properties of this literature are decisive for the present paper.
First, it is overwhelmingly \emph{residential}. The canonical benchmarks
---REDD \cite{kolter2011redd} and UK-DALE---consist of household
appliances whose signatures are sparse, short, and well separated.
Industrial machinery behaves very differently: loads are large,
continuously modulated by process throughput, dominated by induction
motors whose reactive power draw persists at zero mechanical output, and
governed by shift schedules rather than occupant behaviour. Holmegaard
and Kj\ae{}rgaard \cite{holmegaard2016industrial} evaluate established
NILM algorithms on industrial loads and report substantially degraded
performance, attributing it to precisely these differences. Martins
\textit{et al.} \cite{martins2018imdeld} apply a deep generative model
to industrial disaggregation and, in doing so, release the IMDELD
dataset \cite{imdeld2018dataset}---eleven meters recording eight heavy
machines in a Brazilian poultry-feed pellet plant at 1\,Hz across a
five-month window, of which the covered fraction differs by machine
(Section~\ref{sec:method:data}). IMDELD is the benchmark used throughout this
paper; its authors characterise the machines as three-state systems but
supply no per-timestamp state annotation, which is the reason the
present work must proceed without labels and validate its partition
against physics instead.

Second, and more fundamentally, the objective of NILM is
\emph{attribution}: given an aggregate signal, determine which appliance
consumed what. It is a retrospective accounting exercise. NILM does not
forecast, and it does not act. No NILM system, industrial or
residential, closes the loop from an inferred state to an actuation
decision. That step belongs to the two literatures reviewed next.


% ---------------------------------------------------------------------
\subsection{Industrial Energy Optimisation}
\label{sec:rw:industrial}

The manufacturing engineering community has independently established
that non-productive energy dominates the energy budget of discrete-part
production. Gutowski \textit{et al.} \cite{gutowski2006electrical} and
Dahmus and Gutowski \cite{dahmus2004environmental} show that the
load-independent baseline of a machine tool---spindle idling, hydraulics,
coolant pumps, control electronics---frequently exceeds the energy
actually delivered to the workpiece, and that this fraction rises as
automation increases. Kara and Li \cite{kara2011unit} formalise the
relationship as a unit-process model in which specific energy
consumption is inversely related to material removal rate, so that a
lightly loaded or idle machine is intrinsically inefficient. Devoldere
\textit{et al.} \cite{devoldere2007improvement} quantify the resulting
improvement potential and identify idle-mode elimination as the single
largest lever available without redesigning the machine. Duflou
\textit{et al.} \cite{duflou2012towards} synthesise this body of work
into a systems-level roadmap, placing operational measures---switching
equipment off when it is not needed---at the machine and cell levels of
their hierarchy.

On the management side, ISO~50001 \cite{iso50001} provides the
plan--do--check--act framework within which energy performance
indicators are defined and tracked, and ISO~14955-1 \cite{iso14955}
supplies the operating-mode taxonomy that makes standby energy a
measurable quantity in the first place.

The limitation of this literature is methodological rather than
conceptual. Its energy models are typically built from controlled
laboratory measurements of a single machine under prescribed process
parameters, and its recommendations are design guidelines and audit
procedures. It does not address how a facility, given only the
electrical measurements it already collects and no process-parameter
instrumentation, should determine \emph{which} of its machines are idle,
\emph{when}, and for \emph{how long}. That inference problem is left
open.


% ---------------------------------------------------------------------
\subsection{Predictive Shutdown and Decision Support}
\label{sec:rw:decision}

The operations research community has studied the shutdown decision
itself in considerable depth. Mouzon \textit{et al.}
\cite{mouzon2007operational} establish the central trade-off: turning a
machine off saves its idle power but incurs a restart energy penalty and
a restart delay, so shutdown is economical only when the idle period
exceeds a break-even duration determined by the ratio of restart energy
to idle power. They propose dispatching rules and a mathematical
programming formulation that exploit this, and extend the treatment to
joint energy--tardiness objectives \cite{mouzon2008framework}. Chen
\textit{et al.} \cite{chen2013schedule} develop schedule-based operation
policies for production lines, Sun and Li \cite{sun2013opportunity}
estimate real-time control opportunities from buffer state, Li and Sun
\cite{li2016energy} formulate dynamic energy control as a Markov
decision process, Frigerio and Matta \cite{frigerio2015energy} derive
control strategies for machine tools under stochastic job arrivals, and
He \textit{et al.} \cite{he2015energy} embed energy responsiveness in
job-shop machine selection and sequencing.

This literature is mathematically mature, and the break-even
formulation of Mouzon \textit{et al.} is the natural objective for the
decision layer of any standby-elimination system. Its universal
assumption, however, is that the idle interval is \emph{known}---either
because it is dictated by a deterministic production schedule, or
because it is drawn from an analytically tractable arrival process
calibrated by the modeller. The policies are evaluated in simulation or
on analytical queueing models. What is absent is any mechanism for
obtaining the idle-duration input from a machine's own measured
electrical behaviour, which is the only signal available in a retrofit
installation. In other words: this community can decide optimally given
a duration, and the NILM community can infer a state from
measurements, but no work connects the two.


% ---------------------------------------------------------------------
\subsection{Synthesis and Research Gap}
\label{sec:rw:gap}

Table~\ref{tab:related} compares representative works from the four
themes against the capabilities required for an autonomous
standby-elimination system: whether the method infers machine state
without labels, whether it models temporal structure explicitly, whether
it forecasts the \emph{duration} of a future non-productive interval,
whether it converts that forecast into an economically justified control
decision, and whether it is validated on real industrial measurements.

\begin{table*}[!t]
\centering
\caption{Comparison of representative prior work against the
capabilities required for autonomous standby elimination.
\CIRCLE~= provided, \LEFTcircle~= partially provided, \Circle~= not
provided, --~= not applicable.}
\label{tab:related}
\begin{threeparttable}
\small
\setlength{\tabcolsep}{4pt}
\begin{tabular}{@{}lllccccc@{}}
\toprule
\multirow{2}{*}{\textbf{Work}} &
\multirow{2}{*}{\textbf{Method family}} &
\multirow{2}{*}{\textbf{Data domain}} &
\textbf{Labels} &
\textbf{Temporal} &
\textbf{Duration} &
\textbf{Economic} &
\textbf{Real ind.} \\
 & & & \textbf{required} & \textbf{model} & \textbf{forecast} &
 \textbf{decision} & \textbf{data} \\
\midrule
\multicolumn{8}{@{}l}{\textit{Theme A --- machine state identification}}\\
Lloyd \cite{lloyd1982least}          & $k$-means                    & generic                & no  & \Circle    & \Circle & \Circle & \Circle \\
McLachlan \& Peel \cite{mclachlan2000finite} & Gaussian mixture     & generic                & no  & \Circle    & \Circle & \Circle & \Circle \\
Ester \textit{et al.} \cite{ester1996dbscan} & DBSCAN               & generic                & no  & \Circle    & \Circle & \Circle & \Circle \\
Fox \textit{et al.} \cite{fox2011sticky}     & sticky HDP-HMM       & audio                  & no  & \CIRCLE    & \Circle & \Circle & \Circle \\
\addlinespace
\multicolumn{8}{@{}l}{\textit{Theme B --- non-intrusive load monitoring}}\\
Hart \cite{hart1992nonintrusive}     & event/steady-state signature & residential            & yes\tnote{a} & \LEFTcircle & \Circle & \Circle & \Circle \\
Kolter \& Jaakkola \cite{kolter2012afamap} & additive factorial HMM & REDD (residential)     & no  & \CIRCLE    & \Circle & \Circle & \Circle \\
Kelly \& Knottenbelt \cite{kelly2015neuralnilm} & deep NN (dAE, LSTM) & UK-DALE (residential) & yes & \LEFTcircle & \Circle & \Circle & \Circle \\
Zhang \textit{et al.} \cite{zhang2018seq2point} & seq2point CNN     & UK-DALE, REDD          & yes & \LEFTcircle & \Circle & \Circle & \Circle \\
Holmegaard \& Kj\ae{}rgaard \cite{holmegaard2016industrial} & NILM algorithm evaluation & industrial & yes & \LEFTcircle & \Circle & \Circle & \CIRCLE \\
Martins \textit{et al.} \cite{martins2018imdeld} & deep generative disaggregation & IMDELD (industrial) & yes & \LEFTcircle & \Circle & \Circle & \CIRCLE \\
\addlinespace
\multicolumn{8}{@{}l}{\textit{Theme C --- industrial energy optimisation}}\\
Gutowski \textit{et al.} \cite{gutowski2006electrical} & thermodynamic/empirical process model & lab machine tools & -- & \Circle & \Circle & \Circle & \LEFTcircle \\
Kara \& Li \cite{kara2011unit}       & unit-process energy model    & lab machine tools      & --  & \Circle    & \Circle & \Circle & \LEFTcircle \\
Duflou \textit{et al.} \cite{duflou2012towards} & systems-level roadmap & survey             & --  & \Circle    & \Circle & \LEFTcircle & \LEFTcircle \\
\addlinespace
\multicolumn{8}{@{}l}{\textit{Theme D --- predictive shutdown and decision support}}\\
Mouzon \textit{et al.} \cite{mouzon2007operational} & dispatching rules, MILP & simulated single machine & -- & \Circle & \Circle & \CIRCLE & \Circle \\
Chen \textit{et al.} \cite{chen2013schedule} & schedule-based control  & simulated line     & --  & \LEFTcircle & \Circle & \CIRCLE & \Circle \\
Frigerio \& Matta \cite{frigerio2015energy} & stochastic control policy & analytical model  & --  & \CIRCLE    & \LEFTcircle\tnote{b} & \CIRCLE & \Circle \\
He \textit{et al.} \cite{he2015energy} & energy-responsive scheduling & job-shop model      & --  & \Circle    & \Circle & \CIRCLE & \Circle \\
\midrule
\textbf{This work}                   & \textbf{physics-gated GMM--HMM $\rightarrow$ duration forecast $\rightarrow$ optimisation} & \textbf{IMDELD (industrial)} & \textbf{no} & \CIRCLE & \CIRCLE & \CIRCLE & \CIRCLE \\
\bottomrule
\end{tabular}
\begin{tablenotes}[flushleft]\footnotesize
\item[a] Hart's formulation is unsupervised in its clustering step but
requires a manually curated appliance signature library to name the
recovered clusters.
\item[b] Idle duration is modelled as a known stochastic arrival
process calibrated by the analyst, not forecast from measured machine
behaviour.
\end{tablenotes}
\end{threeparttable}
\end{table*}

Reading Table~\ref{tab:related} column-wise makes the gap explicit.
Every capability required for autonomous standby elimination has been
solved somewhere, but the last column of any single row is always
empty:

\begin{itemize}
  \item \textbf{Label-free state inference with temporal structure}
  exists (Theme~A, and \cite{kolter2012afamap} in Theme~B), but is
  demonstrated on generic or residential data and terminates at the
  labelled sequence. Nothing is predicted and nothing is decided.

  \item \textbf{Industrial validation} exists
  (\cite{holmegaard2016industrial,martins2018imdeld}, Theme~B), but the
  task is disaggregation, the methods are supervised, and the output is
  a retrospective attribution of energy rather than an operating state
  usable for control.

  \item \textbf{Quantification of idle-energy waste} exists (Theme~C),
  but rests on controlled laboratory measurement of individual
  processes and yields design guidance, not an online inference
  procedure applicable to a running facility.

  \item \textbf{Optimal shutdown decisions} exist (Theme~D), and the
  break-even formulation is well understood, but every such policy
  takes the idle duration as a \emph{given} input---supplied by a
  deterministic schedule or an analyst-calibrated arrival
  process---and is evaluated in simulation.
\end{itemize}

\medskip
\noindent\fbox{\parbox{0.97\columnwidth}{%
\textbf{Research gap.} No existing work combines
(i)~label-free, temporally coherent inference of industrial machine
states from electrical measurements alone, (ii)~forecasting of the
\emph{duration} of the resulting non-productive intervals, and
(iii)~an economically optimised shutdown decision derived from that
forecast, validated end-to-end on real industrial measurement data.
The state-inference and decision-making capabilities have matured
independently and remain unconnected: the former produces labels no
controller consumes, the latter consumes durations no inference
procedure supplies.}}
\medskip

\noindent
Closing this gap requires more than concatenating existing components.
Three problems arise specifically at the interfaces. First, because no
ground-truth state annotation exists for industrial data, the inferred
partition cannot be validated by conventional supervised metrics, and an
alternative validation basis---independent physical and operational
evidence---must be constructed. Second, an i.i.d.\ clustering produces a
label sequence whose flicker makes duration forecasting meaningless,
so temporal coherence is a prerequisite for the forecasting stage rather
than a refinement of the inference stage. Third, a forecast duration
carries uncertainty that a deterministic break-even threshold ignores;
the decision layer must therefore be posed as a constrained
cost-minimisation problem in which the break-even point emerges as a
derived quantity rather than being supplied as a tuning parameter.
Section~\ref{sec:method} addresses each of these in turn.
